\chapter{Teoretično ozadje}
\label{sec:ozadje}

V tem poglavju predstavimo metodološke pojme, ki so potrebni za razumevanje grafov, grafovskih nevronskih mrež in valence.

\section{Grafi}
\label{subsec:grafi}
% TODO: potem, ko prestavljamo sekcijo drugam, oz. združujemo z drugimi sekcijami in poglavji, poskusi delati tako, da imaš vsaj malo smiselnega teksta med dvemi enumerated deli (sekcijo in podsekcijo). ni pa nujno. ne forsiraj.
\subsection{Osnovna definicija grafa}

Graf definiramo kot
\begin{equation}
    G = (V, E),
\end{equation}
kjer je $V$ množica vozlišč, $E \subseteq V \times V$ pa množica povezav med vozlišči.
Vozlišča praviloma vsebujejo značilke $\mathbf{x}_v$, ki jih zapišemo v matriko
\begin{equation}
    X \in \mathbb{R}^{|V| \times d},
\end{equation}
kjer je $d$ število značilk posameznega vozlišča.
Povezavo med vozliščema $i$ in $j$ označimo z $e_{ij} \in E$.

\subsection{Usmerjeni, uteženi in heterogeni grafi}

V usmerjenih grafih ima povezava določeno izvorno in ponorno vozlišče, zato velja $e_{ij} \ne e_{ji}$.
V neusmerjenih grafih povezave nimajo smeri, zato velja $e_{ij} = e_{ji}$.

Povezave lahko vsebujejo tudi uteži.
Utež povezave med vozliščema $i$ in $j$ označimo z $w_{ij}$.
Za neutežene povezave si lahko predstavljamo, da imajo utež $w_{ij}=1$.
Če imajo povezave več uteži, jih zapišemo kot vektor $\mathbf{e}_{ij}$~\cite{gilmer2017mpnn}.

\subsubsection{Homogeni in heterogeni grafi}
% TODO: prej jih že omenjaš in šele tu opišeš. popravi vrstni red, ampak po pameti.
Homogeni graf je graf, ki ne loči različnih vrst vozlišč in povezav.
Heterogeni graf vsebuje vsaj dve različni vrsti vozlišč ali povezav.
V tej nalogi uporabljamo heterogeno grafovsko predstavitev, saj ločimo več tipov povezav.

\section{Grafovske nevronske mreže}
\label{subsec:gnn}

Grafovske nevronske mreže (GNN) so razred nevronskih mrež, namenjen učenju nad podatki, ki jih lahko predstavimo z grafom~\cite{gilmer2017mpnn,kipf2017gcn,velickovic2018gat}.
Cilj GNN je naučiti se predstavitev vozlišč, povezav ali celotnega grafa tako, da model pri izračunu predstavitve posameznega vozlišča upošteva tudi njegovo lokalno okolico.

Osnovna ideja GNN je \emph{posredovanje sporočil}~\cite{gilmer2017mpnn}.
Vsako vozlišče v posamezni plasti prejme informacije od svojih sosedov, jih agregira in na podlagi zbranih informacij posodobi svojo predstavitev.

Splošno plast posredovanja sporočil lahko zapišemo kot
\begin{equation}
    \mathbf{h}_v^{(k)} =
    \phi^{(k)}
    \left(
        \mathbf{h}_v^{(k-1)},
        \bigoplus_{u \in \mathcal{N}(v)}
        \psi^{(k)}
        \left(
            \mathbf{h}_v^{(k-1)},
            \mathbf{h}_u^{(k-1)},
            \mathbf{e}_{uv}
        \right)
    \right),
\end{equation}
kjer je $\mathcal{N}(v)$ množica sosedov vozlišča $v$, $\psi^{(k)}$ funkcija sporočila, $\oplus$ agregacijska funkcija, $\phi^{(k)}$ funkcija posodobitve, $\mathbf{h}_v^{(k)}$ pa predstavitev vozlišča $v$ v plasti $k$.
Funkcija $\psi^{(k)}$ je lahko linearna preslikava ali večslojni perceptron nad predstavitvama vozlišč in značilkami povezave, $\phi^{(k)}$ pa preprosta vsota prejšnje predstavitve vozlišča in agregiranega sporočila, večslojni perceptron ali rekurentna enota.
Isti funkciji $\psi^{(k)}$ in $\phi^{(k)}$ uporabimo pri vseh vozliščih, zato je plast ob invariantni agregaciji ekvivariantna na permutacijo vozlišč, kar pomeni, da preureditev vozlišč le enako preuredi njihove izhodne predstavitve.
Agregacijska funkcija $\oplus$ mora biti invariantna na vrstni red sosedov.
To lastnost imajo na primer vsota, povprečje in maksimum, zato lahko ista GNN plast deluje na grafih različnih velikosti.
Pri časovnih vrstah je ureditev meritev sicer naravno določena s časom, vendar jo mora model zajeti eksplicitno, na primer s časovnimi povezavami ali časovnimi značilkami, saj se z grafovsko predstavitvijo sicer izgubi.

Pri heterogenem grafu, kjer povezave razdelimo po tipih oziroma relacijah $r \in \mathcal{R}$, namesto ene same množice sosedov obravnavamo množice $\mathcal{N}_r(v)$, ki vsebujejo sosede vozlišča $v$ po relaciji $r$.
Sporočila se najprej agregirajo ločeno znotraj posamezne relacije:
\begin{equation}
    \mathbf{m}_{v,r}^{(k)} =
    \bigoplus_{u \in \mathcal{N}_r(v)}
    \psi_r^{(k)}
    \left(
        \mathbf{h}_u^{(k-1)},
        \mathbf{e}_{uv}
    \right).
\end{equation}
Pri tem je $\psi_r^{(k)}$ funkcija sporočila za relacijo $r$, $\mathbf{m}_{v,r}^{(k)}$ pa agregirano sporočilo, ki ga vozlišče $v$ prejme po tej relaciji v plasti $k$.
Za vsako relacijo lahko uporabimo ločeno linearno preslikavo ali večslojni perceptron $\psi_r^{(k)}$, agregacija znotraj posamezne relacije pa mora biti, tako kot prej, invariantna na vrstni red sosedov.
Relacijska sporočila $\mathbf{m}_{v,r}^{(k)}$ se nato združijo v skupno sporočilo za vozlišče.
Združevanje relacij je lahko vsota ali povprečje, lahko pa vključuje konkatenacijo oziroma naučeno fuzijo relacijskih predstavitev; pri zadnjih dveh mora biti vrstni red relacij vnaprej določen.

\subsection{Klasifikacija na ravni grafa}
\label{subsec:grafovska-klasifikacija}

Če želimo napovedati razred celotnega grafa, lahko nalogo obravnavamo kot klasifikacijo na ravni grafa.
Naj bo $L$ število plasti grafovske nevronske mreže. Na zadnji, $L$-ti plasti dobimo končne predstavitve vozlišč $\mathbf{h}_v^{(L)}$, ki jih združimo v predstavitev celotnega grafa:
\begin{equation}
    \mathbf{h}_G =
    \text{READOUT}
    \left(
        \{\mathbf{h}_v^{(L)} : v \in V\}
    \right).
\end{equation}
Funkcija $\text{READOUT}$ je agregacijska funkcija na ravni grafa.
Pogoste izbire so povprečje, vsota, maksimum ali združevanje s pozornostjo.
Predstavitev grafa $\mathbf{h}_G$ nato uporabimo v klasifikacijski glavi, ki napove razred celotnega grafa.

\subsection{Prekomerno glajenje in kolaps predstavitev}
\label{subsec:oversmoothing}

Pri grafovskih nevronskih mrežah lahko več zaporednih agregacij povzroči, da si predstavitve vozlišč postanejo preveč podobne.
Ta pojav imenujemo prekomerno glajenje.
V praksi se lahko pokaže kot zmanjšanje variance predstavitev, visoka medsebojna podobnost vozliščnih vektorjev in skoraj konstantni izhodi klasifikacijske glave.

Zato smo med učenjem in evalvacijo modelov spremljali statistike predstavitev na različnih delih modela.
Varianca predstavitev nam pokaže, ali model še ohranja razlike med vozlišči in grafi.
Kosinusna podobnost pokaže, ali predstavitve postajajo usmerjene v podobno smer.
Razpon logitov pa pokaže, ali klasifikacijska glava daje dovolj raznolike napovedi ali se približuje skoraj konstantnemu izhodu.
